% !PRINTED 461
\weberpage{461}
\begin{center}
§.\,124.

\gesperrt{The full system of invariants.}
\end{center}

We can now prove that the forms $f, \Delta, C, K$ are a
full system of invariants of the group, i.\,e. that all invariants
of the group can be represented as integral rational functions of these
four. We rely in this on the theorem of
\gesperrt{Bezout} (Bd.\,I, §.\,49), that two curves of the $m$th
and $n$th order which have more than $mn$ points in common
must have a common curve part. Points of contact are
in this to be counted as double or multiple. All intersection points of two
invariant curves either form one connected system, or they
split into systems of connected points, and all points of one
such system are to be counted equally often.

Let $\Phi$ be an invariant of the $m$th order which does not contain the function
$f$ as a factor, and among the intersection points of the curves
$\Phi$ and $f$ let the poles $P_7$ occur $h_1$ times, the poles $P_3$ $h_2$ times,
the poles $\overline{P}_2$ $h_3$ times. Besides these there shall also occur $h_4$
systems of 168 connected points each, of which moreover
(in the case of contact) several systems may coincide into one counted multiply.
Then the number of all intersection points of the two curves is
\[
4m = 24h_1 + 56h_2 + 84h_3 + 168h_4,
\]
and hence
\weq{1}{m = 6h_1 + 14h_2 + 21h_3 + 42h_4,}
in which $h_1, h_2, h_3, h_4$ denote non-negative integers, which
also cannot all vanish. The smallest value that $m$
can have is therefore 6, and an invariant curve of the $6$th
order must pass through the 24 points $P_7$, as we have
already proved for $\Delta$. If $\Delta'$ is a second invariant
of the $6$th order, which must therefore also pass through the points $P_7$,
then we can determine the constant $a$ in $\Delta' - a\Delta$ so
that the invariant curve $\Delta' - a\Delta = 0$ passes through
some $25$th point of $f$. But then, if
$\Delta' - a\Delta$ does not vanish identically,
$\Delta' - a\Delta$ must by \gesperrt{Bezout}'s theorem be divisible
by $f$. The quotient would be an invariant of the second order, which does not
exist; consequently $\Delta' - a\Delta$ must be identically zero.
